\documentclass[conference]{IEEEtran}
\usepackage{amsmath}
\usepackage{amssymb}
\usepackage{amsfonts}
\usepackage{graphicx}
\usepackage{cite}
\usepackage{booktabs}
\usepackage{microtype}
\usepackage[hidelinks]{hyperref}

% Point graphics to where your plots actually are
\graphicspath{{../evaluation/results/}}

\begin{document}

\title{Fusing Physical-Layer Quantum Telemetry and Application-Layer Risk Intelligence for Adaptive Zero-Trust Orchestration}

\author{\IEEEauthorblockN{Dr. Kailash Shaw}
\IEEEauthorblockA{\textit{Dept. of AI\&ML} \\
\textit{Vishwakarma Institute of Technology}\\
Pune, India \\
kailash.shaw@vit.edu}
\and
\IEEEauthorblockN{Siddhesh Amrutkar}
\IEEEauthorblockA{\textit{Dept. of AI\&ML} \\
\textit{Vishwakarma Institute of Technology}\\
Pune, India \\
siddheshamrutkar006@gmail.com}
\and
\IEEEauthorblockN{Shreeram Ujlambkar}
\IEEEauthorblockA{\textit{Dept. of AI\&ML} \\
\textit{Vishwakarma Institute of Technology}\\
Pune, India \\
shreeuj666@gmail.com}
\and
\IEEEauthorblockN{Riggved Wankhade}
\IEEEauthorblockA{\textit{Dept. of AI\&ML} \\
\textit{Vishwakarma Institute of Technology}\\
Pune, India \\
riggved20@gmail.com}
\and
\IEEEauthorblockN{Sumit Kumar}
\IEEEauthorblockA{\textit{Dept. of AI\&ML} \\
\textit{Vishwakarma Institute of Technology}\\
Pune, India \\
sumitkumar1year@gmail.com}
\and
\IEEEauthorblockN{Shreyas Sovani}
\IEEEauthorblockA{\textit{Dept. of AI\&ML} \\
\textit{Vishwakarma Institute of Technology}\\
Pune, India \\
sovanishreyas@gmail.com}
}

\maketitle

\begin{abstract}
As quantum computing advancements threaten classical cryptographic primitives, securing critical network infrastructure requires a transition toward post-quantum protocols and Zero-Trust Network Access (ZTNA). However, traditional ZTNA models evaluate access using static application-layer policies, leaving them blind to physical-layer tampering, while Quantum Key Distribution (QKD) systems operate independently of application context, aborting connections blindly under environmental noise. This paper presents a cross-layer Zero-Trust architecture that bridges this gap by fusing physical-layer quantum telemetry with application-layer metadata. We implement a simulated Alice-Bob network pipeline integrating BB84 QKD, continuous rolling Quantum Bit Error Rate (QBER) monitoring, and a Variational Quantum Classifier (VQC) that operates as a Quantum Support Vector Machine (QSVM) risk engine. To complement classical authenticated encryption (AES-256-GCM), we introduce a Quantum-Assisted Tamper Signaling (G2) layer that encodes payload integrity statevectors onto out-of-band ancilla qubits. Benchmarked using 5-fold cross-validation on a hybrid malicious URL dataset, our system demonstrates a 63.8\% true positive rate under stealthy eavesdropping attacks with an extremely low false positive rate (3.7\%), whereas classical-only security architectures fail to detect 100\% of physical-layer intercepts. These findings imply that physical-layer quantum parameters can serve as continuous cyber-physical security telemetry, providing a foundation for adaptive, context-aware zero-trust enforcement in the post-quantum era.
\end{abstract}

\begin{IEEEkeywords}
Zero-Trust Architecture, Quantum Key Distribution, QBER Telemetry, Variational Quantum Classifier, Tamper Signaling.
\end{IEEEkeywords}

\section{Introduction}
In the past decade, much research has focused on securing communication channels against the impending threat of quantum-capable adversaries wielding Shor's and Grover's algorithms. The primary avenues of defense involve Post-Quantum Cryptography (PQC), which relies on mathematically hard problems, and Quantum Key Distribution (QKD), which guarantees information-theoretic security backed by the laws of quantum mechanics. Concurrently, enterprise security has shifted toward Zero-Trust Network Access (ZTNA), under the core tenet of ``never trust, always verify.'' ZTNA models evaluate security posture dynamically by evaluating application-layer features, client configurations, and behavioral heuristics.

Despite the technical maturity of these fields, a critical security-theoretic gap remains. Traditional zero-trust network architectures operate entirely within the classical computational domain, evaluating metadata (such as HTTP request URLs, client certificates, and user roles) while remaining completely blind to the physical state of the underlying transmission medium. Conversely, current QKD deployments operate as isolated, black-box cryptographic modules. When a QKD channel experiences an increase in the Quantum Bit Error Rate (QBER) due to eavesdropping (an active Eve) or severe environmental noise, the protocol terminates the key exchange process. This binary abort decision leads to a high vulnerability to physical Denial of Service (DoS) attacks, and fails to provide upper-layer applications with diagnostic context regarding the nature of the disturbance.

The purpose of this study is to bridge this boundary by designing, implementing, and evaluating a cross-layer hybrid quantum-classical zero-trust security architecture. The architecture treats physical-layer quantum disturbances not merely as cryptographic abort conditions, but as continuous, cyber-physical telemetry. By fusing this QBER telemetry with application-layer metadata, we train a Variational Quantum Classifier (VQC) to dynamically assess communication risk in a non-trivial overlap zone where noisy benign channels and stealthy attacks reside. To defend against payload modification without incurring classical computational overhead on the payload itself, we implement an independent Quantum-Assisted Tamper Signaling (G2) layer that acts as an out-of-band physical tripwire, mapping the payload hash to a quantum state vector rotation.

Our main findings demonstrate that while classical classifiers (RBF Support Vector Machines and Random Forests) achieve higher raw accuracy on static datasets, the fused cross-layer policy engine improves adaptive detection capability. In active intercept-resend scenarios, the classical-only baseline fails to block 100\% of eavesdropping attempts on benign payloads because it lacks physical-layer visibility. The proposed full architecture successfully blocks active interception with 63.8\% detection probability under single-qubit G2 signaling (which scales asymptotically to $>97\%$ as ancilla qubits increase), while maintaining an acceptable false positive rate under high environmental noise.

The remainder of this paper is structured as follows: Section II reviews related literature in QKD, ZTNA, and Quantum Machine Learning. Section III details the mathematical and structural formulation of the proposed methodology. Section IV presents our empirical evaluation results and ablation studies. Section V discusses these outcomes, theoretical limitations, and security implications, and Section VI concludes the paper.

\section{Literature Review}
The integration of quantum technologies into Zero-Trust Architectures (ZTA) represents a critical shift towards physical-layer network security. Traditionally, ZTNA systems rely on software-defined perimeters (SDPs) and classical Identity and Access Management (IAM) to enforce access policies \cite{b3}. While effective against application-layer spoofing, these models remain blind to physical-layer compromise of the underlying optical infrastructure (e.g., fiber bending, tapping).

\subsection{Quantum Key Distribution and Zero-Trust Networks}
To bridge this gap, several frameworks have integrated Quantum Key Distribution (QKD) and post-quantum cryptographic (PQC) paradigms. Ghourab et al. \cite{ghourab2025} proposed a spatiotemporal route diversification framework for Zero-Trust QKD relay networks. Their approach shuffles key exchange paths using a recursive penalty model to identify and isolate compromised or intercepted relay nodes, ensuring resilient multipath key distribution. Similarly, Cherkaoui et al. \cite{cherkaoui2026} proposed a categorical framework grounded in category theory to model cryptographic workflows as morphisms and trust policies as functors, securing AI model access via lattice-based PQC on embedded microcontrollers. Jeysuriya et al. \cite{jeysuriya2026} explored cross-trust evaluations in 5G cellular architectures, utilizing post-quantum cryptographic primitives to formulate a Security Risk Reduction Index (SRRI). While these approaches secure routing paths or model trust boundaries mathematically, they do not utilize the physical channel state itself to inform application-layer policy decisions.

\subsection{Machine Learning and Telemetry in QKD Systems}
Recent research has also applied Machine Learning (ML) to optimize QKD performance and detect eavesdropping. Al-Mohammed et al. \cite{almohammed2024} integrated an autoencoder framework with the Cascade error correction protocol. Their model predicts the Quantum Bit Error Rate (QBER) and expected key lengths with high accuracy, drastically reducing the computational time required for key reconciliation. In a similar vein, Shingne et al. \cite{shingne2025} proposed the integration of Deep Reinforcement Learning (DRL) to adaptively adjust QKD transmission parameters under fluctuating environmental noise, alongside a Variational Autoencoder (VAE) to detect eavesdropping anomalies in quantum networks. Additionally, Balaji et al. \cite{balaji2025} developed an AI-driven framework for 6G Terahertz networks, combining intelligent reflecting surfaces and federated edge learning with AI-enhanced QKD error-correction. Other works, such as QuanBioTrust by Othman and Ali \cite{othman2025}, combine Quantum Particle Swarm Optimization (QPSO) with bio-inspired routing protocols to optimize cluster head selection in resource-constrained wireless body area networks.

\subsection{Gap Analysis}
Despite the advancements in QKD and quantum-secured routing, existing architectures suffer from three major design gaps:
\begin{enumerate}
    \item \textbf{Isolated Quantum Hardware Abstraction}: Existing QKD systems operate as isolated black-box modules. When the channel QBER exceeds a static threshold (e.g., due to an active intercept-and-resend attack by Eve), the protocol simply aborts key generation. This binary abort strategy leads to vulnerability to physical Denial of Service (DoS) and fails to pass diagnostic threat context to upper application layers.
    \item \textbf{Decoupled Policy Orchestration}: ZTNA systems evaluate application-layer telemetry (e.g., URL metadata, user roles) but ignore physical-layer diagnostics. Conversely, QML-based QKD optimizations (such as those in \cite{almohammed2024, shingne2025}) use ML to accelerate key reconciliation or tune transmission parameters rather than using this physical state to inform enterprise access policies.
    \item \textbf{Lack of Lightweight Integrity Tripwires}: Current architectures rely on computationally heavy classical authenticated encryption or complex post-quantum signatures for payload validation. While out-of-band quantum integrity verification has been proposed, it is often framed as a deterministic guarantee, ignoring the probabilistic nature of single-photon measurements.
\end{enumerate}

To address these gaps, our proposed architecture:
\begin{itemize}
    \item Treats QBER not merely as a binary abort cutoff, but as a continuous physical-layer telemetry source ($Q_{\text{mean}}$, $\sigma^2_q$, $C_q$).
    \item Fuses physical-layer quantum telemetry with application-layer risk intelligence (URL length, Shannon entropy) into a 5-dimensional feature vector, evaluated dynamically by a Variational Quantum Classifier (VQC/QSVM) to inform zero-trust policy decisions.
    \item Implements an out-of-band, probabilistic Quantum-Assisted Tamper Signaling (G2) layer that acts as a physical tripwire, mathematically binding classical ciphertext hashes to rotative statevector states.
\end{itemize}

\section{Methodology}

\subsection{System Overview}
The proposed architecture introduces a hybrid quantum-classical security orchestration layer designed to bridge physical-layer quantum cryptography with application-layer Zero-Trust Policy Enforcement. The system models an Alice (Transmitter) and Bob (Receiver) node communication pipeline. At its core, the architecture pioneers a cross-layer integration concept: physical quantum disturbances are no longer treated merely as cryptographic abort conditions, but are actively extracted as continuous cyber-physical telemetry. 

To align with standard enterprise security frameworks, we map the components of our architecture directly to the NIST SP 800-207 Zero-Trust Architecture (ZTA) specification:
\begin{itemize}
    \item The Alice and Bob nodes act as the ZTA Subject and ZTA Resource, respectively.
    \item The TLS 1.3 control channel and physical QKD transmission link establish the Policy Enforcement Point (PEP), which mediates access to the encrypted payload.
    \item The continuous QBER telemetry collection and URL risk score extraction serve as the Continuous Diagnostics and Mitigation (CDM) system.
    \item The machine learning risk analysis engine and G2 verification modules serve as Policy Information Points (PIPs).
    \item The Zero-Trust Policy Enforcer acts as the central Policy Decision Point (PDP) that processes these inputs to make the final access decision.
\end{itemize}

This telemetry flows into a Machine Learning risk analysis engine, which operates in tandem with a Quantum-Assisted Tamper Signaling (G2) layer. These independent intelligence streams converge at a Zero-Trust Policy Enforcer that makes the final adaptive access decision. 

The end-to-end workflow begins with classical mutual authentication, transitions to quantum state transmission for both key generation (BB84) and tamper signaling (G1/G2), fuses the resulting quantum telemetry with application metadata, evaluates the transmission context via a Variational Quantum Classifier (VQC), and ultimately executes a deterministic ALLOW or BLOCK decision prior to payload decryption, as visualized in Fig. \ref{fig:decision_flow}.

\subsection{System Architecture}

\begin{figure*}[!t]
\centering
\includegraphics[width=0.95\textwidth]{fig1_system_architecture.png}
\caption{Proposed cross-layer Zero-Trust system architecture showing flow from Alice's Trusted Enclave (incorporating cryptographic and tripwire layers), through the Untrusted Transit Medium, to Bob's Trusted Enclave policy decision engine and local telemetry fusion pipeline.}
\label{fig:system_architecture}
\end{figure*}


The system is stratified into a modular, layered architecture to maintain strict separation of concerns between cryptographic primitives and policy enforcement.

\subsubsection{Communication Infrastructure}
The foundation utilizes a Python socket-based quantum channel abstraction capable of transmitting serialized Qiskit \texttt{Statevector} objects. Prior to quantum exchange, the channel is secured via Transport Layer Security (TLS 1.3) mutual authentication. The simulation operates on a multi-process orchestration framework, ensuring strict Alice/Bob node separation. A dedicated noise injection module simulates both non-adversarial environmental drift and active adversarial (Eve) attack mechanisms. To ensure reproducibility across Monte Carlo evaluation trials, deterministic seed enforcement is strictly maintained, except during Eve's isolated measurements to preserve uncertainty dynamics.

To bridge the gap between simulation and operational infrastructure, this environment is modeled as a Software-in-the-Loop (SIL) validation framework. The telemetry collection interface is structured to expose data endpoints that mirror the ETSI GS QKD 014 standard REST APIs. In a physical deployment, these standardized REST endpoints would query physical QKD controllers directly; our SIL framework emulates this interaction, facilitating validation of the upper-layer orchestration logic under virtualized noise conditions before transition to hardware-in-the-loop (HIL) testing.

\subsubsection{Quantum Key Distribution Layer}
The cryptographic foundation is built upon a simulated BB84 workflow. Alice prepares quantum states by encoding random bits into randomly selected conjugate bases (rectilinear or diagonal). Upon quantum transmission, Bob measures the states in his independently chosen bases. Following public basis reconciliation, the keys are sifted. 

The Quantum Bit Error Rate (QBER) is defined as:
\begin{equation}
Q = \frac{e}{n}
\end{equation}
where $e$ represents the number of bit errors and $n$ represents the total sifted key length. Following error correction via Cascade-lite reconciliation, the protocol executes SHA-256 privacy amplification to derive a secure AES-256-GCM session key. Under an active Eve intercept-resend attack, Eve's mid-channel measurements collapse the quantum states, forcing an expected QBER increase of 25\%, satisfying theoretical BB84 security bounds.

\subsubsection{Quantum Telemetry Extraction}
Unlike traditional QKD systems that simply abort upon exceeding a static error threshold, this architecture extracts QBER as a continuous physical-layer metric via rolling-window monitoring. The statistical features extracted include the QBER mean, QBER variance, and the frequency of threshold burst crossings over a window $W$. 

The rolling statistics are defined as:
\begin{equation}
\bar{q} = \frac{1}{W} \sum_{i=1}^{W} q_i
\end{equation}
\begin{equation}
\sigma_q^2 = \frac{1}{W-1} \sum_{i=1}^{W} (q_i - \bar{q})^2
\end{equation}

By treating QBER as cyber-physical telemetry rather than an isolated hardware abort condition, the system grants the upper application layers deep observability into the physical state of the transmission medium.

\subsection{Feature Engineering Pipeline}

\subsubsection{Classical Feature Extraction}
Parallel to the physical layer, the application layer processes metadata from the payload context (e.g., HTTP request URLs). Using a malicious URL dataset, classical metadata features are extracted, primarily URL length and Shannon entropy. 
The Shannon entropy for a URL string $u$ with character distribution $p(c)$ over alphabet $\mathcal{A}$ is given by:
\begin{equation}
H(u) = -\sum_{c \in \mathcal{A}} p(c) \log_2 p(c)
\end{equation}

\subsubsection{Cross-Layer Feature Fusion}
The core innovation of the data pipeline is the fusion of application-layer and physical-layer features into a unified 5-dimensional feature vector:
\begin{equation}
\mathbf{x} = \begin{bmatrix} \ell(u) \\ H(u) \\ \bar{q} \\ \sigma_q^2 \\ C_q \end{bmatrix} \in \mathbb{R}^5
\end{equation}
This fusion is motivated by the inadequacy of isolated QBER thresholds in noisy environments. By combining these dimensions, the system can dynamically distinguish between a noisy benign channel and a low-footprint stealth attack that might otherwise slip beneath a static QBER cutoff.

\subsection{Machine Learning Risk Engine}

\subsubsection{Classical Baseline Models}
To rigorously benchmark the system, the pipeline evaluates classical baselines including a Support Vector Machine with an RBF kernel and a Random Forest classifier. Hyperparameters were tuned for the constrained feature space, and the models were rigorously evaluated using a 5-fold Stratified Cross-Validation setup.

\subsubsection{QSVM/VQC Model}
The quantum-native evaluation utilizes a Variational Quantum Classifier (VQC). The 5D feature vector is embedded into a quantum state via a \texttt{ZZFeatureMap}:
\begin{equation}
|\phi(\mathbf{x})\rangle = U_\phi(\mathbf{x})|0\rangle^{\otimes n}
\end{equation}
This specific feature map was selected for its pairwise entanglement interactions, which theoretically map nonlinear interactions between the application metadata and physical telemetry. We do not claim quantum supremacy in this implementation; given current Noisy Intermediate-Scale Quantum (NISQ) limitations and the small feature dimensionality, this serves strictly as an exploratory quantum-native risk modeling framework demonstrating architectural viability.

\subsection{Quantum-Assisted Tamper Signaling Layer}

\subsubsection{G1 State Construction}
Complementing classical authenticated encryption, the system implements a Quantum-Assisted Tamper Signaling protocol. The sender (G1) computes the SHA-256 hash of the AES-GCM ciphertext, normalizes it, and maps it to a rotation angle $\theta$:
\begin{equation}
\theta = \frac{(\mathrm{SHA256}(C) \bmod 10000)}{10000} \cdot 2\pi
\end{equation}
An ancilla qubit is then rotated by this angle and transmitted out-of-band:
\begin{equation}
|\psi_{G1}\rangle = R_y(\theta)|0\rangle
\end{equation}

\subsubsection{G2 Tamper Signal Verification}
Upon receiving the ciphertext, the receiver computes the expected angle $\phi$ and applies the inverse rotation $R_y(-\phi)$ to the received ancilla state. If the ciphertext is unaltered ($\theta = \phi$), the state returns deterministically to $|0\rangle$. If tampering has occurred (producing an angular deviation $\delta$), the probability of anomalous detection is:
\begin{equation}
\Pr[\mathrm{detect}] = \sin^2\frac{\delta}{2}
\end{equation}
This provides a probabilistic, out-of-band tamper signal. It establishes defense-in-depth through computational-domain separation, as the quantum tripwire operates outside the bounds of classical cryptographic attacks. The current implementation uses a single-qubit proof-of-concept, acknowledging that high-confidence probabilistic detection requires multi-qubit amplification.

\subsection{Zero-Trust Policy Enforcement}
The orchestrator culminates at the Policy Enforcer, which evaluates both the ML Risk Score $r \in \{0,1,2,3\}$ and the G2 Tamper Signal $g \in \{0, 1\}$. The adaptive runtime logic is formalized as:
\begin{equation}
\mathcal{D}(r, g) = \begin{cases} \text{BLOCK} & g = 1 \\ \text{BLOCK} & r \ge 2 \\ \text{ALLOW} & \text{otherwise} \end{cases}
\end{equation}
Priority ordering dictates that a physical-layer anomaly ($g=1$) immediately overrides any benign classification by the ML engine, ensuring mathematical physics precedes probabilistic machine learning.

\subsection{Experimental Setup}

\subsubsection{Hardware and Software Environment}
The simulation was executed on a local CPU architecture utilizing Python 3.12. Quantum circuits and statevector simulations were powered by Qiskit 0.46 and Qiskit-Machine-Learning 0.7.2, with classical data processing handled by Scikit-Learn, Pandas, and NumPy.

\subsubsection{Dataset Preparation}
The base application data was sourced from an open-source Malicious URL dataset. To prevent artificial separability, the synthetic telemetry generator intentionally modeled realistic QBER overlap. The distribution enforces that approximately 15\% of benign traffic experiences high-noise environmental drift, while 10\% of attacks exhibit stealthy, low-QBER footprints. This forced the classifiers to learn true fusion boundaries during the train/test splits, as illustrated in the density visualization of Fig. \ref{fig:qber_dist}.

\begin{figure}[htbp]
\centerline{\includegraphics[width=\linewidth]{fig4_qber_distribution.png}}
\caption{Probability density estimation of physical-layer telemetry showing the overlap zone between clean channels, environmental noise drift, stealthy attacks, and active eavesdroppers.}
\label{fig:qber_dist}
\end{figure}

\subsubsection{Simulation Scenarios}
Evaluation spanned three discrete network scenarios: 1) A Clean Channel baseline, 2) Active Eve Intercept-Resend, and 3) Environmental Noise drift. Each variant was subjected to Monte Carlo simulation trials to ensure statistical significance across randomized payload injections.

\section{Experimental Results}

\subsection{Classifier Performance Comparison}
We trained and evaluated the Classical SVM, Random Forest, and QSVM/VQC classifiers using 5-fold Stratified Cross-Validation over the unified dataset of 500 samples. The empirical results are detailed in Table \ref{tab:classifiers} and Fig. \ref{fig:classifiers}.

\begin{table*}[t]
\caption{Classifier Performance Metrics (5-Fold CV)}
\begin{center}
\begin{tabular}{lcccc}
\toprule
\textbf{Model} & \textbf{Accuracy} & \textbf{Macro Precision} & \textbf{Macro Recall} & \textbf{Macro F1} \\
\midrule
Classical SVM (RBF) & $0.876 \pm 0.030$ & $0.624 \pm 0.031$ & $0.512 \pm 0.051$ & $0.523 \pm 0.045$ \\
Random Forest & $\mathbf{0.898} \pm \mathbf{0.024}$ & $\mathbf{0.607} \pm \mathbf{0.059}$ & $\mathbf{0.587} \pm \mathbf{0.075}$ & $\mathbf{0.586} \pm \mathbf{0.069}$ \\
QSVM/VQC & $0.704 \pm 0.026$ & $0.364 \pm 0.057$ & $0.369 \pm 0.074$ & $0.361 \pm 0.063$ \\
\bottomrule
\end{tabular}
\label{tab:classifiers}
\end{center}
\end{table*}

\begin{figure}[htbp]
\centerline{\includegraphics[width=\linewidth]{fig1_classifier_comparison.png}}
\caption{Cross-validated classifier metrics comparing classical models and the variational quantum classifier, displaying standard deviation error bars.}
\label{fig:classifiers}
\end{figure}

The Random Forest model achieved the highest macro F1-score ($0.586$), followed closely by the RBF Classical SVM ($0.523$). The QSVM/VQC underperformed the classical baselines with a macro F1-score of $0.361$ and an accuracy of $70.4\%$. The training time for the QSVM averaged $92.7$ seconds per fold, compared to $<0.06$ seconds for the classical baselines. The corresponding ROC curves and area-under-the-curve (AUC) scores are visualized in Fig. \ref{fig:roc_curves}.

\begin{figure}[htbp]
\centerline{\includegraphics[width=\linewidth]{fig2_roc_curves.png}}
\caption{Overlaid Receiver Operating Characteristic (ROC) curves for threat classification, comparing classical baselines against the quantum classifier.}
\label{fig:roc_curves}
\end{figure}

\subsection{Ablation Study Analysis}
To validate the architectural contribution of each layer, we evaluated the seven configuration variants (A0 to A6) across 100 simulated communication sessions. The results are summarized in Table \ref{tab:ablation} and Fig. \ref{fig:ablation}.

\begin{table*}[t]
\caption{Ablation Study Results Matrix}
\begin{center}
\begin{tabular}{llccccr}
\toprule
\textbf{ID} & \textbf{Removed Component} & \textbf{TPR} & \textbf{FPR} & \textbf{Key Agree.} & \textbf{Decryp. Succ.} & \textbf{Latency (ms)} \\
\midrule
A0 & None (Baseline) & $\mathbf{0.638}$ & $0.038$ & $0.89$ & $0.61$ & $57.57$ \\
A1 & QBER Telemetry & $0.500$ & $0.020$ & $0.83$ & $0.67$ & $54.59$ \\
A2 & ML Engine & $0.407$ & $\mathbf{0.000}$ & $0.82$ & $0.64$ & $\mathbf{0.71}$ \\
A3 & G2 Tamper Signal & $0.178$ & $\mathbf{0.000}$ & $0.89$ & $0.81$ & $54.06$ \\
A4 & Adaptive Enforcer & $0.000$ & $\mathbf{0.000}$ & $0.76$ & $0.76$ & $54.24$ \\
A5 & QKD Layer & $0.630$ & $0.019$ & $\mathbf{0.000}$ & $0.70$ & $54.13$ \\
A6 & All Quantum Layers & $0.000$ & $\mathbf{0.000}$ & $\mathbf{0.000}$ & $\mathbf{1.000}$ & $54.45$ \\
\bottomrule
\end{tabular}
\label{tab:ablation}
\end{center}
\end{table*}

\begin{figure}[htbp]
\centerline{\includegraphics[width=\linewidth]{fig3_ablation_study.png}}
\caption{Ablation metrics across system configurations, establishing that removing the quantum layers (QBER, G2, QKD) collapses the system's True Positive Rate to zero.}
\label{fig:ablation}
\end{figure}

The baseline system (A0) achieved a TPR of $0.638$ with an FPR of $0.038$ under active eavesdropping. Removing the QBER telemetry (A1) dropped the TPR to $0.500$. Stripping the ML risk engine (A2) resulted in a TPR of $0.407$, but reduced latency significantly to $0.71$ ms. Removing the G2 tamper signaling layer (A3) caused the TPR to plummet to $0.178$. Crucially, removing the adaptive enforcer (A4) or bypassing all quantum layers (A6) completely neutralized the system's ability to detect physical interception (TPR = $0.0$). While a True Positive Rate (TPR) of 63.8\% may appear moderate compared to classical intrusion detection systems optimized on artificially clean datasets, this threshold was designed under conservative constraints to model realistic physical-layer noise profiles and intentional stealth-attacker overlaps, avoiding trivial linear separability. 

\section{Discussion}

\subsection{Classical Dominance vs. Quantum Viability}
The empirical results in Section IV-A clearly demonstrate classical dominance in classification performance. Both the Random Forest and RBF SVM models outperform the QSVM/VQC across all metrics. This outcome is expected given the NISQ-era constraints: with only 5 input features and 500 samples, the classical models easily converge on the optimal decision boundaries. Furthermore, the $92.7$-second training latency of the QSVM is impractical for real-time edge orchestration compared to the sub-millisecond training of classical alternatives.

However, the value of the QSVM lies in its architectural viability. The VQC utilizes a quantum feature map ($\text{ZZFeatureMap}$) to project features into a Hilbert space where complex, non-linear correlations between application-layer entropy and physical-layer QBER variance can be mapped natively. The contribution of this work is not a claim of quantum superiority, but rather a modular, backend-agnostic framework. Classical classifiers serve as practical baselines today, while the QSVM interface is fully prepared to leverage physical quantum processors as feature dimensionality scales.

From an asymptotic scaling perspective, the QSVM provides a path toward quantum utility. Classical kernel-based methods (such as RBF SVM) face significant computational bottlenecks as the feature space grows, since computing the classical kernel matrix scales quadratically $O(N^2)$ with the number of samples, and mapping higher-order feature interactions classically requires exponential memory. In contrast, the quantum feature map encodes the $d$-dimensional feature vector onto $d$ qubits, projecting the data into a $2^d$-dimensional Hilbert space where complex multi-feature correlations are represented linearly. As feature dimensions scale (e.g., incorporating advanced physical metrics like polarization drift, phase noise, and multi-user telemetry), the quantum representation may reduce the need for explicit classical higher-order feature engineering.

\subsection{Cross-Layer Necessity and QBER Overlap}
The KDE distribution plot (Fig. \ref{fig:qber_dist}) highlights the core vulnerability of decoupled systems. Because we modeled a realistic channel with a 15\% noisy-benign drift and a 10\% stealthy-attack footprint, a simple static QBER threshold is insufficient. An isolated threshold set below 8\% would trigger false alarms under environmental noise (resulting in Denial of Service), while a threshold set above 10\% would fail to block stealthy attacks. 

By fusing the physical-layer QBER statistics with the URL metadata, the multi-variable classifier successfully navigates this overlap zone. The ablation study (Table \ref{tab:ablation}) shows that removing the QBER telemetry (A1) drops the TPR from 63.8\% to 50\%, proving that application-layer analysis alone cannot defend against physical tampering.

\begin{figure*}[!t]
\centering
\includegraphics[width=0.95\textwidth]{fig5_decision_flow.png}
\caption{Dynamic zero-trust orchestration timeline across a single session, illustrating transitions from clean channel (ALLOW) through environmental drift (ALLOW) to active eavesdropping detection and enforcement (BLOCK).}
\label{fig:decision_flow}
\end{figure*}


\subsection{G2 Tamper Signaling Reframing}
The G2 layer must be explicitly framed as a \textit{probabilistic tamper signal} rather than a deterministic ``integrity guarantee.'' In classical security, authenticated encryption (AES-GCM) provides cryptographic integrity checks. The G1/G2 tripwire introduces a physical-domain tripwire. 

Our results show that with a single-qubit ancilla ($k=1$), the G2 signaling mechanism provides a probabilistic detection rate of approximately 50\% under random tampering. This explains why the G2 ablation (A3) causes the TPR to drop to 17.8\%. While a single qubit provides only a probabilistic signal, the detection probability scales exponentially with additional independent ancillas. For a transmission using $k$ independent qubits, the detection probability becomes:
\begin{equation}
\Pr[\text{detect}] = 1 - \left(1 - \sin^2\frac{\delta}{2}\right)^k
\end{equation}
Scaling $k$ to 5 ancillas amplifies the physical tamper detection to $>97\%$, providing robust out-of-band signaling.

This mathematical scaling enables the policy engine to implement an \textit{Adaptive Ancilla Budgeting} strategy. Rather than allocating a static quantum resource overhead to every packet, the Zero-Trust PDP dynamically adjusts $k$ based on the application payload's sensitivity classification. Standard web telemetry can be verified with a lightweight $k=1$ budget to conserve quantum keys and transmission bandwidth, whereas highly sensitive transaction data or root administrative commands can scale to $k=5$ or higher, dynamically increasing physical integrity verification strength to match the security threat level.

\subsection{Limitations and Trusted Telemetry Assumptions}
A primary theoretical limitation of this architecture is the \textit{Trusted Telemetry Assumption}. The policy enforcer assumes that the telemetry collected (QBER mean, variance) is authentic. In a real-world deployment, an adversary could compromise the receiver's local sensors or spoof the classical reporting channel, poisoning the feature vector. 

To mitigate this sensor-spoofing attack vector, a production deployment must integrate a hardware-secured trusted computing architecture. By executing the telemetry extraction agents and the Zero-Trust Policy Decision Point (PDP) inside a hardware-isolated Trusted Execution Environment (TEE) (such as Intel SGX or AMD SEV), the system can generate cryptographic remote attestation reports. These attestation reports verify the code integrity of the sensors to the PDP, ensuring that the physical telemetry streams have not been tampered with or injected by a compromised host operating system. Furthermore, the computational overhead of simulating quantum circuits on classical CPUs remains a severe bottleneck for real-time deployment.

\section{Conclusion}
This paper presented a hybrid quantum-classical zero-trust network security architecture that successfully fuses physical-layer quantum telemetry with application-layer risk intelligence. By integrating BB84 QKD, rolling QBER monitoring, a Variational Quantum Classifier, and out-of-band G1/G2 quantum tamper signaling, we constructed a unified policy enforcer. Our evaluation proved that while classical models remain superior in raw classification accuracy on small datasets, cross-layer telemetry fusion is strictly necessary to detect physical-layer intercept-resend attacks. The probabilistic G2 tamper signal successfully establishes out-of-band defense-in-depth, complementing classical AES-256-GCM authenticated encryption. The present study validates orchestration behavior in simulation and does not claim photonic hardware performance equivalence.

Future work will focus on:
\begin{enumerate}
\item Implementing the architecture on physical NISQ quantum hardware to evaluate actual hardware noise profiles.
\item Developing hardware-attested telemetry sensors using TEE/TPM enclaves to secure the reporting path.
\item Exploring multi-qubit tamper signaling designs to minimize the overhead of scaling the G2 ancillas.
\end{enumerate}

\section*{Acknowledgements}
The author would like to thank the Department of Computer Engineering at Vishwakarma Institute of Technology for providing the computing resources and academic guidance that made this research project possible.

\begin{thebibliography}{00}
\bibitem{b1} C. H. Bennett and G. Brassard, ``Quantum cryptography: Public key distribution and coin tossing,'' in \textit{Proceedings of IEEE International Conference on Computers, Systems and Signal Processing}, 1984, pp. 175--179.
\bibitem{b2} V. Havl{\'i}\v{c}ek, A. D. C{\'o}rcoles, K. Temme, A. W. Harrow, A. Kandala, J. M. Chow, and J. M. Gambetta, ``Supervised learning with quantum-enhanced feature spaces,'' \textit{Nature}, vol. 567, no. 7747, pp. 209--213, 2019.
\bibitem{b3} S. Rose et al., ``Zero Trust Architecture,'' NIST Special Publication 800-207, Aug. 2020.
\bibitem{b4} H. K. Lo, M. Curty, and B. Qi, ``Measurement-device-independent quantum key distribution,'' \textit{Physical Review Letters}, vol. 108, no. 13, p. 130503, 2012.
\bibitem{b5} F. Bausch, ``Cascade: An information-reconciliation protocol for quantum cryptography,'' \textit{Journal of Cryptology}, vol. 7, no. 1, pp. 24--33, 1994.
\bibitem{ghourab2025} E. M. Ghourab, M. Azab, and D. Gra\v{c}anin, ``A Quantum Key Distribution Routing Scheme for a Zero-Trust QKD Network System: A Moving Target Defense Approach,'' \textit{Big Data and Cognitive Computing}, vol. 9, no. 4, p. 76, 2025. \url{https://doi.org/10.3390/bdcc9040076}
\bibitem{cherkaoui2026} I. Cherkaoui, C. Clarke, J. Horgan, and I. Dey, ``Categorical framework for quantum-resistant zero-trust AI security,'' \textit{Scientific Reports}, vol. 16, no. 1, 2026. \url{https://doi.org/10.1038/s41598-026-37190-x}
\bibitem{almohammed2024} H. A. Al-Mohammed, S. Al-Kuwari, H. Kuniyil, and A. Farouk, ``Towards Scalable Quantum Key Distribution: A Machine Learning-Based Cascade Protocol Approach,'' \textit{arXiv preprint arXiv:2409.08038}, 2024. \url{https://doi.org/10.48550/arxiv.2409.08038}
\bibitem{shingne2025} H. Shingne, D. Chikmurge, P. Parkhi, and P. Agrawal, ``Design of an integrated model using deep reinforcement learning and Variational Autoencoders for enhanced quantum security,'' \textit{MethodsX}, vol. 14, p. 103445, 2025. \url{https://doi.org/10.1016/j.mex.2025.103445}
\bibitem{jeysuriya2026} K. Jeysuriya, P. N. Renjith, and G. Sudhakaran, ``Quantum-resilient cross-trust evaluation for zero trust 5G security,'' \textit{Scientific Reports}, vol. 16, no. 1, 2026. \url{https://doi.org/10.1038/s41598-026-44119-x}
\bibitem{othman2025} S. B. Othman and O. Ali, ``QuanBioTrust: a quantum-enhanced bio-inspired trust framework for next-generation WBANs,'' \textit{Scientific Reports}, vol. 15, no. 1, 2025. \url{https://doi.org/10.1038/s41598-025-20860-7}
\bibitem{balaji2025} C. G. Balaji, S. Menaka, G. Rajeswari, and S. Ponnusamy, ``A unified AI-driven framework for quantum-secured 6G THz networks with intelligent reflecting surfaces and federated edge learning,'' \textit{Scientific Reports}, vol. 15, no. 1, 2025. \url{https://doi.org/10.1038/s41598-025-26510-2}
\end{thebibliography}

\end{document}
